\chapter{Installation}

\section{Prerequisites}
\index{prerequisites}\index{Raspberry Pi}\index{touch display}\index{hardware requirements}%
The system is designed around the following hardware and software.

\subsection{Hardware}
\begin{itemize}
  \item \textbf{Delta X2 delta robot}, connected over USB serial.
  \item \textbf{Conveyor belt} with a G-code speed controller, connected over
    USB serial.
  \item \textbf{USB camera} mounted above the belt (default capture
    resolution $1280 \times 720$).
  \item \textbf{Target computer: Raspberry Pi 4 or 5} running a 64-bit Linux
    (ARM64/aarch64).
  \item \textbf{Operator display: official Raspberry Pi 7\,'' Touch Display}
    ($800 \times 480$, DSI, capacitive touch). The user interface is designed
    touch-first for this resolution: finger-sized controls, no scrolling on
    the main operator view, and the emergency-stop button always remains the
    largest target on screen.
\end{itemize}

\subsection{Software}
\begin{itemize}
  \item Linux only (see \emph{Supported platform} below); no Windows or macOS
    support is intended.
  \item Rust toolchain $\geq$ 1.85 (the crate uses the 2024 edition).
  \item The system packages listed in the next section (OpenCV, libclang,
    libudev).
\end{itemize}

Development and testing do not require the robot hardware: an x86\_64 Ubuntu
machine can run the full application with simulated hardware
(\cfg{cargo run -{}- -{}-mock}). The production build path for the Raspberry
Pi (cross-compilation from the development machine versus building on the Pi
itself) is not yet fixed; until it is, the repository deliberately carries no
deployment tooling.

\section{Supported platform}
Linux is the supported platform; Ubuntu 24.04 or newer is recommended.
Other distributions work if the packages below are available.

\section{System packages}
\index{installation}\index{system packages}\index{OpenCV}\index{libclang}\index{libudev}%
The build needs a C/C++ toolchain, OpenCV development headers, libclang
(for generating the OpenCV Rust bindings) and libudev (for serial-port
discovery):

\begin{lstlisting}[language=bash]
sudo apt update && sudo apt install -y \
    clang libclang-dev llvm-dev \
    libudev-dev libopencv-dev pkg-config
\end{lstlisting}

\begin{table}[h]
  \centering
  \begin{tabularx}{\textwidth}{lD}
    \toprule
    \textbf{Package} & \textbf{Why it is needed} \\
    \midrule
    \cfg{clang}, \cfg{libclang-dev}, \cfg{llvm-dev} &
      The \cfg{opencv} crate generates Rust bindings from the C++ headers
      at build time using libclang. \\
    \cfg{libopencv-dev} & OpenCV itself (video capture, image processing). \\
    \cfg{libudev-dev} & Required by the \cfg{serialport} crate for USB/serial
      device discovery. \\
    \cfg{pkg-config} & Lets the build scripts locate the installed libraries. \\
    \bottomrule
  \end{tabularx}
  \caption{System package roles}
\end{table}

\section{Rust toolchain}
\index{Rust!toolchain}\index{rustup}%
Install the stable Rust toolchain via \href{https://rustup.rs}{rustup}:
\begin{lstlisting}[language=bash]
curl --proto '=https' --tlsv1.2 -sSf https://sh.rustup.rs | sh
\end{lstlisting}

\section{Serial port permissions}
\index{serial port!permissions}\index{dialout group@\texttt{dialout} group}%
On Linux the serial devices belong to the \cfg{dialout} group:
\begin{lstlisting}[language=bash]
sudo usermod -a -G dialout $USER
# Log out and back in for the change to take effect.
\end{lstlisting}

\section{Building and first run}
\index{building}\index{cargo@\texttt{cargo}}\index{mock mode}%
\begin{lstlisting}[language=bash]
cd deltax2sort
cargo build            # compile (also builds the Slint UI)
cargo test             # run the unit test suite
cargo run -- --mock    # start with simulated hardware
cargo run              # start against the real hardware
\end{lstlisting}

\begin{tipbox}
  Always bring the system up in mock mode first
  (\cfg{cargo run -- --mock}) after installing or changing the
  configuration. Mock mode exercises the full control flow --- UI,
  orchestrator, queueing, E-stop --- without moving any hardware.
\end{tipbox}

\section{Troubleshooting}
\index{troubleshooting}\index{LIBCLANG PATH@\texttt{LIBCLANG\_PATH}}%
\begin{description}[style=nextline]
  \item[\cfg{clang-sys} cannot find \cfg{libclang}]
    Install \cfg{libclang-dev}. If only a versioned library such as
    \cfg{libclang-21.so.1} exists, create a symlink named
    \cfg{libclang.so} (the repository keeps one under \cfg{libs/}) and
    point the build at it:
    \begin{lstlisting}[language=bash]
ln -sf /usr/lib/x86_64-linux-gnu/libclang-21.so.1 libs/libclang.so
LIBCLANG_PATH=$PWD/libs cargo build
    \end{lstlisting}
  \item[\cfg{libudev-sys} build failure]
    \cfg{libudev-dev} is missing; install it and rebuild.
  \item[OpenCV build failure]
    Ensure \cfg{libopencv-dev} and \cfg{pkg-config} are installed and that
    \cfg{pkg-config --modversion opencv4} prints a version.
  \item[\emph{Permission denied} opening \cfg{/dev/ttyUSB*}]
    Your user is not in the \cfg{dialout} group yet, or you have not
    logged out and back in.
  \item[Robot fails the \cfg{IsDelta} handshake]
    Check the port assignment (robot vs.\ conveyor may be swapped), the
    baud rate, and that no other program holds the port open. See
    Appendix~\ref{app:gcode} for the handshake protocol.
\end{description}
